\documentclass[10pt,a4paper]{amsart}
\usepackage[margin=19mm]{geometry}
\usepackage{amsmath,amssymb,mathtools}
\usepackage[scr=rsfs,cal=euler]{mathalfa}
\usepackage{xfrac}
\usepackage[unicode,hidelinks,bookmarks=false]{hyperref}
\newcommand{\ie}{i.e.\ }
\newcommand{\cf}{cf.\ }
\newcommand{\resp}{respectively\ }
\newcommand{\Subsection}{\subsection}
\providecommand{\nobreakdash}{\nobreak\hbox{-}}
\setlength{\emergencystretch}{2em}
\multlinegap0pt
\usepackage{bm}
\usepackage{bbm}
\usepackage{dsfont}
\usepackage{xspace}
\usepackage{cancel}
\usepackage{mathtools}
\usepackage{amsthm}
\makeatletter
\@ifundefined{thm}{\newtheorem{thm}{Theorem}}{}
\makeatother
\newcommand{\R}{{\mathbb{R}}}
\title[Tree-tubings and the combinatorics of resurgent Dyson-Schwin]{Tree-tubings and the combinatorics of resurgent Dyson-Schwinger equations}
\author{Michael Borinsky, Gerald V. Dunne, Karen Yeats}
\date{Source: arXiv:2408.15883v3 (CC BY 4.0)}
\begin{document}
\maketitle

\noindent\emph{Source and licence.} Tree-tubings and the combinatorics of resurgent Dyson-Schwinger equations. Version arXiv:2408.15883v3 (CC BY 4.0), distributed under the Creative Commons Attribution 4.0 International licence, as recorded in the arXiv licence metadata for this exact version and shown on the abstract page https://arxiv.org/abs/2408.15883v3. Full rights notice: licenses/arxiv-2408.15883-CC-BY-4.0.txt.\par\smallskip

\noindent\textit{Editorial scope.} Counted unit: the Thm whose source label is \texttt{thm fixed alpha beta} in the article (source0.tex lines 769--775, source label \texttt{thm fixed alpha beta}), delivered together with the article's own complete original proof of it (source0.tex lines 776--793, one \texttt{proof} environment of 18 lines). No mathematics is written, repaired or re-derived: statement and proof are the article's own text line for line. Required context retained uncounted: ode (lines 712--712); beta\_eq (lines 765--765). Every definition, display and label of those lines is retained unchanged. Omitted from this package and not redistributed: 1--711, 713--764, 766--768, 794--1201. Nothing inside a retained excerpt is omitted or altered: every retained body is delivered exactly as the article's lines are written, so no omission receipt of any kind accompanies this package. Citations: the retained excerpts carry no citation command at all, so no bibliography entry of the article is transcribed and no reference is redistributed. Reference adaptation: none was needed and none was made; every reference inside the delivered unit resolves inside it, so no unresolved reference or citation is produced. Printed slips kept literal: no printed word, symbol, hypothesis or number of the counted statement or of its proof was altered. Layout and class: the article's own class is \texttt{amsart} with packages including geometry, graphicx, hyperref, fontenc, inputenc, lmodern, amssymb, amsmath, xcolor; the wrapper is the lane's pinned \texttt{amsart} editorial wrapper at 10pt on A4, which restates the theorem-like environments the retained text needs and adds the article's own macro definitions that the retained text needs (\texttt{\textbackslash R}), with extra wrapper packages (bm, bbm, dsfont, xspace, cancel, mathtools). The delivered numbering is the unit's own. Assets: no retained span contains a figure environment, a graphics-inclusion command, a PDF-page inclusion, a table or a TikZ picture; no image file is copied into this package, the package declares no asset dependency and no image file is redistributed with it. Licence: version arXiv:2408.15883v3 of this article is distributed under the Creative Commons Attribution 4.0 International licence, as recorded in the arXiv licence metadata for this exact version and shown on the abstract page https://arxiv.org/abs/2408.15883v3. A full rights notice is delivered at licenses/arxiv-2408.15883-CC-BY-4.0.txt. Characters: every retained excerpt is pure ASCII, so the delivered engine, XeLaTeX with its default Latin Modern text font, reports no missing character.
\begin{align} \label{ode} \gamma_j(x) \,-\, \gamma_{j-1}(x) \,= \,\gamma(x) \left( 2 x \frac{\partial}{\partial x} \,-\,1 \right) \gamma_j(x) \text{ for } 1 \leq j \leq m\,, \end{align}

\begin{align} \label{beta_eq} 0 &= - d_{q-1} - a_q \left( c_m - (2 \beta+1) c_q + 2 \alpha^{-1} c_q \sum_{k=1}^m a_k\right). \end{align}

\begin{thm}
\label{thm fixed alpha beta}
If all the $a_1,\ldots,a_m$ are different and non-vanishing and the unique set of solution power series $\gamma_1,\ldots,\gamma_m \in \R[[x]]$ of the equation system in eq.~\eqref{ode} are elements of $\R^\alpha_\beta[[x]]$ for some $\alpha,\beta \in \R$, then there is an integer $q$, $1 \leq q \leq m$ such that $(\alpha,\beta) = (\alpha_q,\beta_q)$, where 
\begin{align} \label{eq:alphas} \alpha_q &\,=\, 2 a_q \quad\text{ and }\\
\label{eq:betas} \beta_q &\,=\, \frac{1}{2 a_q}\sum_{\substack{i=1\\ i\neq q}}^m a_i \,+\, \frac12 \prod_{\substack{i=1\\ i\neq q}}^{m}\frac{a_q}{a_q-a_i}\,. \end{align}
Here, the empty sum is $0$, and the empty product is $1$, as usual.
\end{thm}

\begin{proof}
Substitution of the values for $d_{r-1}$ and $c_m$ 
into Eq.~\eqref{beta_eq}  results in 
\begin{align*} \beta &= \frac{1}{2 a_q} \left( -a_q + \sum_{j=1}^q a_j \prod_{\substack{i=j\\ i\neq q}}^{m} \frac{a_q}{a_q-a_i} + \sum_{i=1}^m a_i \right). \end{align*}
It remains to be proven that 
\[
     \sum_{j=1}^q a_j \left( \prod_{\substack{i=j\\ i\neq q}}^{m}\frac{a_q}{a_q-a_i} \right) = a_q \left( \prod_{\substack{i=1\\ i\neq q}}^{m}\frac{a_q}{a_q-a_i} \right),
\]
which is equivalent to 
    \[
    \sum_{j=1}^q a_j \prod_{i=1}^{j-1}\frac{a_q-a_i}{a_q} = a_q.
    \]
We prove this equality by breaking the sum into two pieces
    \[
    \sum_{j=1}^q a_j \prod_{i=1}^{j-1}\frac{a_q-a_i}{a_q} = \underbrace{\sum_{j=1}^{q-1} a_j \prod_{i=1}^{j-1}\left(1 - \frac{a_i}{a_q}\right)}_A + \underbrace{a_q \prod_{i=1}^{q-1}\left(1 - \frac{a_i}{a_q}\right)}_B.
    \]
    Every term in $A$ has the form $(-1)^ta_j/a_q^t$ times a product of $t$ $a_i$s with $i<j$ for some $t\geq 0$. Here, $t$ is the number of factors of the product from which the $a_i/a_q$ was taken to obtain this term.  Furthermore, every such term appears exactly once in $A$.  For $B$, the term where $1$ is taken from each factor when expanding the product is $a_q$.  For every other term in $B$, the explicit $a_q$ can be canceled, so every such term has the form $(-1)^{t+1}/a_q^t$ times a product of $t+1$ $a_i$s for some $t\geq 0$, where here $t+1$ is the number of times the $a_i/a_q$ was taken when expanding out the product, and again every such term appears exactly once. Organizing these terms by the largest index $j$ appearing in the product of $a_i$s we see these terms are exactly the terms obtained from $A$ but with opposite signs.  Therefore the terms of $A$ cancel with all terms of $B$ except for the term $a_q$, giving the desired result.
\end{proof}

\end{document}
